Das Volumen des Pneus bleibt praktisch gleich, der Druck ist also proportional zur absoluten Temperatur. Mit $T_1 = \TEiF = \TEi$ und $T_2 = \TZwF = \TZw$ ist
\begin{align}
 p_2 &= \pZwF = \pEi\cdot\frac{\TZw}{\TEi} = \pZw \approx \result{\pZwbP{0}{2}}
\end{align}
Bemerkung: Der Reifendruck wird meist als Überdruck $p_\text{ü}$ gegenüber der Umgebung angegeben. Ist \pEiO\ ein Überdruck, so muss man im Gasgesetz den absoluten Druck $p_\text{ü}+p_0$ (mit $p_0 = \ncpo$) verwenden, und der Überdruck am Ende ist
\begin{align}
 p_\text{ü,2} &= \pUeF = \pUe \approx \pUebP{0}{2}
\end{align}
